In this section you create the storage and delivery layer for the web
interface. An S3 bucket holds the files; CloudFront delivers them to
users.

One thing to note: the bucket will stay \textbf{completely private}.
Users never reach S3 directly, only through CloudFront. The mechanism
that makes this work is called \textbf{Origin Access Control}.

\vspace{0.5em}\noindent\textbf{Step 1. Create the frontend
bucket}\par

\begin{enumerate}
\def\labelenumi{\arabic{enumi}.}
\tightlist
\item
  Open the
  \href{https://s3.console.aws.amazon.com/s3/home?region=ap-southeast-1}{Amazon
  S3 console}
\item
  Choose \textbf{Create bucket}
\item
  Fill in:

  \begin{itemize}
  \tightlist
  \item
    \textbf{Bucket name}:
    \texttt{fcj-recsys-frontend-\textless{}your-name\textgreater{}}
    (bucket names are globally unique, so add a suffix)
  \item
    \textbf{AWS Region}: Asia Pacific (Singapore)
    \texttt{ap-southeast-1}
  \item
    \textbf{Block Public Access}: \textbf{leave everything enabled}
  \end{itemize}
\item
  Choose \textbf{Create bucket}
\end{enumerate}

\begin{figure}
\centering
\pandocbounded{\includegraphics[keepaspectratio,alt={Create bucket}]{5-Workshop/5.3/create-bucket.png}}
\caption{Create bucket}
\end{figure}

Older guides tell you to enable Static website hosting and make the
bucket public. That approach is no longer recommended because it exposes
the bucket to the Internet. Here we keep the bucket closed and let
CloudFront be the only door.

\vspace{0.5em}\noindent\textbf{Step 2. Build the frontend}\par

On your machine, go into the \texttt{frontend} folder and build:

\begin{Shaded}
\begin{Highlighting}[]
\BuiltInTok{cd}\NormalTok{ frontend}
\ExtensionTok{npm}\NormalTok{ install}
\ExtensionTok{npm}\NormalTok{ run build}
\end{Highlighting}
\end{Shaded}

This produces a \texttt{dist} folder containing all static files.

Do not worry about the API address yet. After you create API Gateway in
section 5.5, you will come back, set \texttt{VITE\_API\_BASE\_URL} and
rebuild.

\vspace{0.5em}\noindent\textbf{Step 3. Upload to S3}\par

\begin{Shaded}
\begin{Highlighting}[]
\ExtensionTok{aws}\NormalTok{ s3 sync dist/ s3://fcj{-}recsys{-}frontend{-}}\OperatorTok{\textless{}}\NormalTok{your{-}name}\OperatorTok{\textgreater{}}\NormalTok{/ }\AttributeTok{{-}{-}delete}
\end{Highlighting}
\end{Shaded}

The \texttt{-\/-delete} flag removes old files no longer present in the
new build, so leftovers do not pile up across deployments.

Check the console: the bucket should now contain \texttt{index.html} and
an \texttt{assets} folder.

\vspace{0.5em}\noindent\textbf{Step 4. Create the CloudFront
distribution}\par

\begin{enumerate}
\def\labelenumi{\arabic{enumi}.}
\tightlist
\item
  Open the
  \href{https://console.aws.amazon.com/cloudfront/v4/home}{CloudFront
  console}
\item
  Choose \textbf{Create distribution}
\item
  Under \textbf{Origin}:

  \begin{itemize}
  \tightlist
  \item
    \textbf{Origin domain}: pick the bucket you just created
  \item
    \textbf{Origin access}: choose \textbf{Origin access control
    settings (recommended)}
  \item
    Choose \textbf{Create new OAC}, keep the defaults, click
    \textbf{Create}
  \end{itemize}
\item
  Under \textbf{Default cache behavior}:

  \begin{itemize}
  \tightlist
  \item
    \textbf{Viewer protocol policy}: \textbf{Redirect HTTP to HTTPS}
  \item
    \textbf{Allowed HTTP methods}: \texttt{GET,\ HEAD}
  \end{itemize}
\item
  Under \textbf{Settings}:

  \begin{itemize}
  \tightlist
  \item
    \textbf{Default root object}: \texttt{index.html}
  \end{itemize}
\item
  Choose \textbf{Create distribution}
\end{enumerate}

\begin{figure}
\centering
\pandocbounded{\includegraphics[keepaspectratio,alt={Create distribution}]{5-Workshop/5.3/create-distribution.png}}
\caption{Create distribution}
\end{figure}

\textbf{Advanced Security:} At this CloudFront creation step, you could
also apply \textbf{AWS WAF} (Web Application Firewall) to block Botnets
and DDoS attacks directed at your web interface.

\vspace{0.5em}\noindent\textbf{Step 5. Grant CloudFront read access to the
bucket}\par

After creation, CloudFront shows a banner prompting you to update the
bucket policy. Choose \textbf{Copy policy}, then:

\begin{enumerate}
\def\labelenumi{\arabic{enumi}.}
\tightlist
\item
  Go to the
  \href{https://s3.console.aws.amazon.com/s3/home?region=ap-southeast-1}{S3
  console} and open the bucket
\item
  Open the \textbf{Permissions} tab
\item
  Under \textbf{Bucket policy}, choose \textbf{Edit}, paste the copied
  policy, then \textbf{Save changes}
\end{enumerate}

This policy allows only that one distribution to read the bucket, nobody
else.

\vspace{0.5em}\noindent\textbf{Step 6. Handle 403 errors for the single-page
app}\par

The React frontend uses client-side routing. When a user opens
\texttt{/cart} directly, CloudFront looks for a file named \texttt{cart}
in S3, does not find it and returns an error. Those errors need to be
redirected to \texttt{index.html}.

\begin{enumerate}
\def\labelenumi{\arabic{enumi}.}
\tightlist
\item
  In the distribution, open the \textbf{Error pages} tab
\item
  Choose \textbf{Create custom error response} and set:

  \begin{itemize}
  \tightlist
  \item
    \textbf{HTTP error code}: \texttt{403\ Forbidden}
  \item
    \textbf{Customize error response}: Yes
  \item
    \textbf{Response page path}: \texttt{/index.html}
  \item
    \textbf{HTTP Response code}: \texttt{200\ OK}
  \end{itemize}
\item
  Repeat once more for error code \texttt{404\ Not\ Found}
\end{enumerate}

\vspace{0.5em}\noindent\textbf{Verify}\par

Wait around 5 minutes for the distribution status to become
\textbf{Deployed}, then open the address shown under
\textbf{Distribution domain name}, in the form
\texttt{dxxxxxxxxx.cloudfront.net}.

The interface should render. Product data will be missing because the
API does not exist yet, which is expected.

\textbf{Remember this one, it will save you hours.} Every time you
upload a new build to S3 you \textbf{must} create an invalidation,
otherwise CloudFront keeps serving the cached old version and you will
think your code is broken.

\begin{Shaded}
\begin{Highlighting}[]
\ExtensionTok{aws}\NormalTok{ cloudfront create{-}invalidation }\AttributeTok{{-}{-}distribution{-}id} \OperatorTok{\textless{}}\NormalTok{YOUR{-}ID}\OperatorTok{\textgreater{}}\NormalTok{ {-}{-}paths }\StringTok{"/*"}
\end{Highlighting}
\end{Shaded}
